\begin{flushright}
  \fechaclase{10 de noviembre de 2026}
\end{flushright}

\subsection{Terminal Sliding Mode Control}

En el SMC convencional se utiliza una superficie $s=e_2+me_1$ con convergencia exponencial $\dot{e}_1=me_1$, que salvo por el tiempo de alcance, la convergencia no ocurre en tiempo finito.\\

El control terminal modifica la superficie para que la dinámica no sea lineal y llegue al origen en tiempo finito, y tendrá un tiempo de alcance $Tr$ finito y un tiempo de deslizamiento $Ts$ finito. EL tiempo total de convergencia es 

\begin{equation*}
    T_{total}=Tr-Ts
\end{equation*}

sea el sistema no lineal SISO\\

\begin{equation*}
    \begin{aligned}
        \dot{x}_1=x_2\\
        \dot{x}_2=f(x,t)+u\\
        y=x_2
    \end{aligned}
\end{equation*}

con errores de seguimiento

\begin{equation*}
    \begin{aligned}
        e_1=&r-y=r-x_1\\
        e_2=&\dot{r}-\dot{y}=\dot{r}-\dot{x}_1=\dot{r}-\dot{x}_2
    \end{aligned}
\end{equation*}

derivando

\begin{equation*}
    \begin{aligned}
        \dot{e}_1=&\dot{r}-\dot{x}_1=\dot{r}-x_2=e_2\\
        \dot{e}_2=&\ddot{r}-\dot{x}_2=\ddot{r}-f(x,t)-u
    \end{aligned}
\end{equation*}

se introduce el operador potencial signada (sigpow)

\begin{equation*}
    <x>^P=|x|^P sign(x)
\end{equation*}

\begin{equation*}
    <f(x)>^P \begin{cases}
        |f(x)^P| & f(x)>0\\
        0 & f(x)=0\\
        -|f(x)^P| & f(x)<0
    \end{cases}
\end{equation*}

cuya derivada es

\begin{equation*}
    \frac{d}{dt}<f(t)^P=P|f(t)|^{P-1}\dot{f}(t)
\end{equation*}

tiene la propiedad

\begin{equation*}
    |Z|^{\gamma-1}<Z>^{2-\gamma}=Z \quad , \quad 1<\gamma<2
\end{equation*}

proponemos una dinámica no lineal para el error

\begin{equation*}
    \dot{e}_1=-m<e_1>^\alpha
\end{equation*}

con $m>0$ y $0<\alpha<1$, que implica

\begin{equation*}
    e_2=-m<e_1>^\alpha
\end{equation*}

Entonces proponemos la superfice de deslizamiento

\begin{equation*}
    S=e_2+m<e_1>^*
\end{equation*}

y cuando $s=0$ $\dot{e}_1=e_2=-m<e_1>^\alpha$ separando variables y definiento $Z=|e_1|$

\begin{equation*}
    \begin{aligned}
        \dot{Z}&=-mZ^\alpha\\
        Z^{-\alpha}dZ&=-mdt
    \end{aligned}
\end{equation*}


Integrando desde $Z(Tr)$ hasta $0$

\begin{equation*}
    Ts|=\frac{|e_1(Tr)|^{1-\alpha}}{m(1-\alpha)}
\end{equation*}

mostrando convergencia en timepo finitio. Se deriva la superficie:

\begin{equation*}
    \begin{aligned}
        \dot{s}&=\dot{e}_2+m\alpha|e_1|^{\alpha-1}\dot{e}_1\\
            &=\dot{e}_2+m\alpha|e_1|^{\alpha-1}e_2\\
            &=m\alpha|e_1|^{\alpha-1}e_2+\ddot{r}-f(x,t)-u
    \end{aligned}
\end{equation*}

proponemos la función de Lyapunov

\begin{equation*}
    \begin{aligned}
        V=&\frac{1}{2}s^2\\
        \dot{V}=&s\dot{s} =s[m\alpha|e_1|^{\alpha-1}e_2+\ddot{r}-f(x,t)-u]
    \end{aligned}
\end{equation*}

bajo la suposición $|\ddot{r}-f(x,t)|\leq F$ para algún F positivo

\begin{equation*}
    \dot{V} \leq s[m\alpha |e_1|^{\alpha-1}e_2+F-u]
\end{equation*}

proponemos $u=m\alpha |e_1|^{\alpha-1}e_2+Ksign(s)$

\begin{equation*}
    \begin{aligned}
        \dot{V} \leq& F|s|-K|s|\\
            \leq & -(K-F)|s|
    \end{aligned}
\end{equation*}

mostrnado convergencia del error en tiempo finito con la ley de contorl

\begin{equation*}
    \begin{aligned}
        &s=e_2+m<e_1>^\alpha\\
        &u= m\alpha|e_1|^{\alpha-1}e_2 + Ksign(s)\\
        &<e_1>^\alpha=|e_1|^\alpha sign(e_1)
    \end{aligned}
\end{equation*}